\label{chap:83-07}

La connexion se construit à partir de la classification orbitale complète et de
son premier relèvement coinductif. Le développement suivant fixe les trois
histoires initiales, les frontières orientées et la chaîne
\(w_6\to w_{30}\to R_{36}\) avant d'introduire l'holonomie nonadique.

\begin{HMTScientificOwnerSubordinated}
\input{ampliaciones_sucesoras_20260824/parte_i_ii/owners/U008_orbitas_elevacion_r36}
\end{HMTScientificOwnerSubordinated}

% Unidad U009. Conexión nonádica conjunta; redacción autónoma.

\section{Connexion nonadique conjointe, compression et mémoire terminale}
\label{p12-83-c7-s1:chap:conexion-nonadica-conjunta}

Les neuf phases qui apparaissent après \(R_{36}\) ne sont pas neuf preuves
indépendantes. Ce sont les neuf composantes locales d'une unique connexion
sur des états enrichis. Leur composition autour d'un tour définit
une holonomie. Le retour de l'étiquette de phase conserve la continuité du
préfixe, de la retenue, de la frontière, de l'ombre de Witt et de la mémoire.

\subsection{Transport de l'état enrichi par l'holonomie nonadique et mise à jour de la mémoire}

Au niveau \(K\), nous écrivons
\[
 \widehat x_K=
 (B_K,C_K,p_K,c_K,\Sigma_K,W_K,g(K),q_{{\rm mem},K}),
\]
où :
\begin{enumerate}
 \item \(B_K=(u_K^\pi,u_K^e,u_K^\varphi)\) est le bloc visible de trois
       mots de six trits ;
 \item \(C_K\) est le cylindre ternaire compatible ;
 \item \(p_K\) est le préfixe complet ;
 \item \(c_K\) est la retenue accumulée ;
 \item \(\Sigma_K\) est la signature conjointe de frontière ;
 \item \(W_K\) est l'ombre d'incidence de Witt ;
 \item \(g(K)=1+((K-1)\bmod9)\) est l'étiquette publique de phase ;
 \item \(q_{{\rm mem},K}\in\mathbb Z_{\geq0}\) est la mémoire verticale non
       bornée.
\end{enumerate}
La coordonnée dodécaphasique est dérivée :
\[
 \beta_K=[q_{{\rm mem},K}]_{12}.
\]
On ne remplace pas \(q_{\rm mem}\) par \(\beta\), car la réduction modulo
douze perdrait le nombre de tours.

\subsection{Relations locales et holonomie}

Pour chaque \(K\), soit
\[
 \mathcal R_K\subseteq
 \widehat X_K\times\widehat X_{K+1}
\]
la relation d'extension compatible entre états complets. La composition
relationnelle d'un tour est
\begin{equation}
 \Gamma_{9,K}
 :=\mathcal R_{K+8}\circ\cdots\circ\mathcal R_{K+1}\circ\mathcal R_K
 :\widehat X_K\longrightarrow\widehat X_{K+9}.
\label{p12-83-c7-s1:eq:gamma9-composicion-local}
\end{equation}

\begin{definition}[Holonomie nonadique]
\label{p12-83-c7-s1:def:holonomia-nonadica}
L'holonomie d'un tour est
\[
 \Gamma_9:=\operatorname{Hol}_{C_9}(\nabla_{\rm TPK}),
\]
dont la réalisation au niveau \(K\) est
\(\Gamma_{9,K}\). Les neuf relations \(\mathcal R_K,\ldots,
\mathcal R_{K+8}\) sont des cartes de cette unique connexion.
\end{definition}

\begin{theorem}[Retour de phase sans réinitialisation]
\label{p12-83-c7-s1:thm:retorno-nonadico-sin-reinicio}
Si \((\widehat x_K,\widehat x_{K+9})\in\Gamma_{9,K}\), alors
\[
 g(K+9)=g(K),
 \qquad
 q_{{\rm mem},K+9}=q_{{\rm mem},K}+1,
\]
et
\[
 \beta_{K+9}=\beta_K+1\pmod{12}.
\]
En général, \(\widehat x_{K+9}\neq\widehat x_K\).
\end{theorem}

\begin{proof}
La première identité provient de la définition de \(g\). Un tour complet
incrémente par définition le compteur de tours non borné. Réduire cette
identité modulo douze donne la troisième. Les autres coordonnées sont
transportées par la composition \cref{p12-83-c7-s1:eq:gamma9-composicion-local} ; aucune
loi du TPK ne les réinitialise. L'égalité de phase n'implique donc pas
l'égalité de l'état complet.
\end{proof}

\subsection{Classification exhaustive des \(9\) phases du Topological Phase Kernel}

Les blocs visibles et les recensements locaux du premier tour sont
\[
\begin{array}{c|rrrrr|ccc}
K&N_\pi&N_e&N_\varphi&N_{\rm loc}&N_{\rm ext}
 &u^\pi&u^e&u^\varphi\\\hline
1&378&422&349&3&2&012222&121221&112202\\
2&238&368&388&2&5&010211&102011&121020\\
3&184&284&173&5&4&002111&012222&010210\\
4&207&207&122&4&1&110221&102011&010200\\
5&39&151&151&1&1&222220&021222&102011\\
6&110&110&69&1&3&111201&201202&221021\\
7&81&29&81&3&3&212121&222210&200221\\
8&59&59&59&3&2&200121&212212&110110\\
9&43&43&43&2&1&100100&020112&010122
\end{array}.
\]

\begin{theorem}[Fonctions des neuf composantes]
\label{p12-83-c7-s1:thm:funciones-nueve-componentes}
Dans le tour précédent, les neuf composantes remplissent les fonctions
suivantes :
\begin{enumerate}
 \item ouverture locale de trois orientations et première résolution ;
 \item séparation binaire avec un maximum local de cinq extensions ;
 \item multiplicité locale maximale de cinq états ;
 \item contraction de quatre états locaux en une extension ;
 \item première saturation forte, avec une seule continuation ;
 \item présent local unique et ouverture triple du futur ;
 \item fibre spéculaire triple ;
 \item synchronisation des trois recensements en \(59\) ;
 \item clôture orientée : deux états locaux d'entrée et une relation qui,
       après l'horizon supplémentaire, conserve une prolongation.
\end{enumerate}
\end{theorem}

\begin{proof}
Chaque affirmation se lit dans les recensements \(N_{\rm loc},N_{\rm ext}\), dans
l'égalité ou l'inégalité de \(N_\pi,N_e,N_\varphi\) et dans la relation
d'extension explicite aux neuvième, dixième et onzième phases qui sera démontrée
dans l'unité suivante. Le tableau énumère le tour entier ; aucun
sous-ensemble de phases n'a été choisi.
\end{proof}

La cinquième composante coïncide avec la première frontière conjointe
\[
 R_{36}=(222220|021222|102011)^T.
\]
La neuvième composante contient
\[
 B_9=(100100|020112|010122).
\]
La présence de deux états locaux à la neuvième phase et de trois extensions
à partir de l'état orienté relève de couches différentes ; ce ne sont pas des recensements
contradictoires.

L'holonomie agit sur l'état enrichi complet. Avant de publier
son quotient modal, on construit donc l'involution qui échange
les fibres de \(\pi\) et \(e\), l'axe fixe de \(\varphi\), les invariants
ternaires et la résolution locale \(3\to1\) déterminée par le deuxième
horizon.

\begin{HMTScientificOwnerSubordinated}
\input{ampliaciones_sucesoras_20260824/parte_i_ii/owners/U010_espejo_seleccion_tres_a_uno}
\end{HMTScientificOwnerSubordinated}

\subsection{Transfert conjoint de la fibre de 468 émissions}

La factorisation du catalogue est
\[
 \mathcal U_6^{\rm cat}\cong
 \mathfrak S_+\times\mathfrak S_\times,
 \qquad
 |\mathfrak S_+|=18,\qquad |\mathfrak S_\times|=26.
\]
Pour une fonction \(f\) sur le catalogue, soient
\[
 (E_+f)(s,e)=\frac1{18}\sum_{s'\in\mathfrak S_+}f(s',e),
\]
et
\[
 (E_\times f)(s,e)=\frac1{26}
 \sum_{e'\in\mathfrak S_\times}f(s,e').
\]
Les projecteurs \(E_+\) et \(E_\times\) commutent.

Soit
\[
 \tau=(+1,-1,0,+1,-1,0,+1,-1,0),
 \qquad
 P_{+1}=E_+,quad P_{-1}=E_\times,quad P_0=I.
\]
Dans \(\mathcal U_6^{\rm cat}\times C_9\), nous définissons le transfert de phase
par
\begin{equation}
 \mathcal K_{\rm ph}\bigl((u,r),(v,[r+1]_9)\bigr)
 :=P_{\tau(r)}(u,v),
 \label{p12-83-c7-s1:eq:u009-kph-explicita}
\end{equation}
et le rendons nul pour les autres couples de phases. Ainsi, la permutation cyclique
et l'opérateur appliqué à chacun des neuf pas sont déterminés.

\begin{proposition}[Renouvellement conjoint après un tour]
\label{p12-83-c7-s1:prop:transferencia-conjunta-nueve}
\[
 \mathcal K_{\rm ph}^{\,9}=E_+E_\times.
\]
En particulier, sur une masse ponctuelle,
\[
 (E_+E_\times)(u,v)=\frac1{18\cdot26}=\frac1{468}.
\]
\end{proposition}

\begin{proof}
Le produit ordonné sur un tour est
\[
 P_{\tau(8)}P_{\tau(7)}\cdots P_{\tau(0)}
 =I E_\times E_+ I E_\times E_+ I E_\times E_+
 =E_+^3E_\times^3I^3=E_+E_\times.
\]
La deuxième égalité utilise la commutativité et la dernière l'idempotence.
Le deuxième résultat est la distribution uniforme sur le produit des
deux signatures.
\end{proof}

Cette identité exprime un renouvellement conjoint, non neuf filtrations
indépendantes.

\subsection{Décomposition de l'holonomie \(U_{\Gamma_9}\) en compression visible \(T\) et composante orthogonale \(\eta\) : \(I-T^*T=\eta^*\eta\)}

Soient \(\mathcal H_{\rm vis}\) et \(\mathcal H_{\rm full}\) des espaces de
Hilbert pour la réalisation visible et la réalisation complète. Soit
\[
 J:\mathcal H_{\rm vis}\longrightarrow\mathcal H_{\rm full}
\]
une inclusion isométrique et
\(U_{\Gamma_9}\) une réalisation isométrique de l'holonomie. Nous définissons
\[
 T:=J^*U_{\Gamma_9}J,
 \qquad
 \eta:=(I-JJ^*)U_{\Gamma_9}J.
\]

\begin{theorem}[Décomposition compression--expression]
\label{p12-83-c7-s1:thm:conexion-compresion-expresion}
On a
\[
 \boxed{U_{\Gamma_9}J=JT+\eta,}
\]
\[
 J^*\eta=0,
 \qquad
 \boxed{I-T^*T=\eta^*\eta.}
\]
\end{theorem}

\begin{proof}
Nous décomposons \(U_{\Gamma_9}J\) au moyen des projections orthogonales
\(JJ^*\) et \(I-JJ^*\) :
\[
 U_{\Gamma_9}J
 =JJ^*U_{\Gamma_9}J+(I-JJ^*)U_{\Gamma_9}J
 =JT+\eta.
\]
L'orthogonalité donne \(J^*\eta=0\). Puisque \(J\) et
\(U_{\Gamma_9}\) sont des isométries,
\[
 I=J^*U_{\Gamma_9}^*U_{\Gamma_9}J
 =T^*T+\eta^*\eta.
\]
Le réarrangement prouve la dernière identité.
\end{proof}

\(T\) est la compression visible ; \(\eta\) est l'expression complémentaire.
Aucune des deux ne doit être appelée contraction stricte. Un opérateur distinct
est réservé à cette fonction.

\subsection{Contraction stricte et télescopage terminal}

Soit \(V_9\) une isométrie et soit \(q\) un paramètre avec \(0<q<1\). Nous définissons
\[
 M_9=qV_9,
 \qquad
 D_9=(I-M_9^*M_9)^{1/2}.
\]
Alors \(\|M_9\|=q<1\) et
\[
 I-M_9^*M_9=D_9^*D_9.
\]

\begin{theorem}[Identité de Stein avec terme terminal]
\label{p12-83-c7-s1:thm:stein-telescopia-terminal}
Pour tout \(L\geq1\),
\[
 \boxed{
 I=
 \sum_{r=0}^{L-1}M_9^{*r}D_9^*D_9M_9^r
 +M_9^{*L}M_9^L.}
\]
Plus généralement, si \(S_0\) satisfait
\(S_0-M_9^*S_0M_9=D_9^*D_9\), alors
\[
 \boxed{
 S_0=
 \sum_{r=0}^{L-1}M_9^{*r}D_9^*D_9M_9^r
 +M_9^{*L}S_0M_9^L.}
\]
\end{theorem}

\begin{proof}
Nous multiplions l'identité de défaut à gauche par \(M_9^{*r}\) et à
droite par \(M_9^r\). En sommant pour \(r=0,\ldots,L-1\), les termes
intérieurs se télescopent et les termes initial et terminal demeurent. La même preuve
s'applique à l'équation générale de Stein.
\end{proof}

Le dernier terme ne s'efface pas à profondeur finie. Sa disparition dans une
limite exige une preuve de convergence ; elle ne fait pas partie de l'identité
algébrique.

\subsection{Cocycle de retenue et puissances relationnelles}

Pour des flèches composables \(\gamma_1,\gamma_2\), la retenue satisfait
\[
 \operatorname{Carr}(\gamma_2\gamma_1)
 =\operatorname{Carr}(\gamma_2)H(\gamma_1)
 +Y(\gamma_2)\operatorname{Carr}(\gamma_1).
\]
La formule conserve la contribution de la première étape après
l'avoir transportée par la seconde.

Les tours de plus grande longueur s'écrivent avec des indices décalés :
\begin{align*}
 \Gamma_{27,K}
 &=\Gamma_{9,K+18}\circ\Gamma_{9,K+9}\circ\Gamma_{9,K},\\
 \Gamma_{54,K}
 &=\Gamma_{9,K+45}\circ\cdots\circ\Gamma_{9,K},\\
 \Gamma_{108,K}
 &=\Gamma_{9,K+99}\circ\cdots\circ\Gamma_{9,K}.
\end{align*}
Ce n'est qu'après avoir déclaré un opérateur global sur l'union graduée que
cette notation peut être abrégée en \(\Gamma_9^3,\Gamma_9^6,\Gamma_9^{12}\).

\subsection{État nonadique complet : phase, retenue, frontière, route et mémoire après la prolongation}

La connexion d'un tour est enregistrée par le quintuplet
\[
 \mathfrak N_9=
 \bigl(\Gamma_9,T,\eta,\operatorname{Carr},S_{\rm term}\bigr),
\]
où \(S_{\rm term}=M_9^{*L}S_0M_9^L\) à une profondeur finie
\(L\). Un foncteur d'application qui utilise la connexion doit déclarer sa
compatibilité avec les cinq champs. Commuter avec la seule phase visible ne
constitue pas une naturalité par rapport à l'état nonadique complet.

\subsection{Obstructions de l'holonomie \(\Gamma_9\) : retour \(9\to1\), indices de composition et conservation de la phase, de la retenue, de la frontière, de la route et de la mémoire}

\begin{criterion}[Critères de réfutation de la connexion]
La construction est réfutée si :
\begin{enumerate}
 \item les neuf phases sont présentées comme des filtres indépendants ;
 \item le retour \(9\to1\) réinitialise \(q_{\rm mem}\), le préfixe ou la retenue ;
 \item \(T\) est appelé contraction stricte ;
 \item une identité de la
       \cref{p12-83-c7-s1:thm:conexion-compresion-expresion} échoue ;
 \item le terme terminal est éliminé à profondeur finie ;
 \item \(\Gamma_{27,K}\) est composé sans décaler les indices ;
 \item un foncteur d'application ignore l'un des cinq champs de
       \(\mathfrak N_9\).
\end{enumerate}
\end{criterion}

Les développements complets qui suivent reconstruisent d'abord l'émetteur fini,
sa factorisation et ses microfibres ; ils déploient ensuite l'involution
spéculaire, les trois extensions de \(B_9\), l'unicité à deux horizons et
l'inverse de Bockstein--Hensel qui conserve la mémoire pendant le retour de phase.

\begin{HMTScientificOwnerSubordinated}
\input{sections/hmt/03b_alfabeto_factorizacion}
\end{HMTScientificOwnerSubordinated}

\begin{HMTScientificOwnerSubordinated}
\input{sections/hmt/14_numero_heisenberg_y_d108}
\end{HMTScientificOwnerSubordinated}
